% Caption for chronic_oral_pct_bars.pdf  (Fig6, fig:dosing_comparison_bars)
% Panels: A) IS, B) HA, C) pCS. Clinical CI on A and C (HA has no clinical data).
% Bars: Total / Bound / Unbound are equal %Δ at constant fu_p (same bar height).
% y-axis shared 0–65%. Clinical CI only on A and B.

\begin{figure}[htbp]
\centering
\includegraphics[width=\textwidth]{Figures/Fig6.pdf}
\caption{%
\textbf{Predicted last-cycle PBUT plasma accumulation under three oral
furosemide dosing regimens compared with the clinical reference.}
Panels \textbf{A)} (indoxyl sulfate, IS), \textbf{B)} (hippuric acid, HA),
and \textbf{C)} ($p$-cresyl sulfate, pCS) show the percentage change
($\Delta C_{ss}$, \%) in time-averaged plasma concentration relative to the
no-drug CKD-5 baseline
($C_{ss,\mathrm{IS}}=28.16$\,mg/L,
$C_{ss,\mathrm{HA}}=142.86$\,mg/L,
$C_{ss,\mathrm{pCS}}=24.99$\,mg/L).
For each regimen three bars are drawn: total (black), bound (dark grey), and
unbound (light grey). At constant $f_{u,p}$, these three $\Delta C_{ss}$ values
are identical, so the bars coincide in height.
Regimens (left to right): TID 120\,mg/day ($3\times40$\,mg),
QD 240\,mg/day ($1\times240$\,mg), and TID 240\,mg/day ($3\times80$\,mg).
In panels A and C the grey band is the 95\% confidence interval of the
clinical reference
($+31.9$\,\% IS [13.5--53.3\,\%];
$+29.3$\,\% pCS [6.1--57.8\,\%];
Costes-Albrespic et~al.\ \cite{costes2025furosemide},
patients prescribed $>120$\,mg/day furosemide equivalent);
the dashed line is the clinical mean.
Panel B has no clinical HA comparator (annotation on the panel).
All ODE predictions use the same CKD-5 parameters
($K_i = 0.35$\,mg/L calibrated to \cite{costes2025furosemide};
$F = 0.60$ \cite{FDA_furosemide_label};
$k_a = 1.5$\,h$^{-1}$ \cite{Hammarlund1984}).
Panels A--C share the $y$-axis range $0$--$65$\,\% so IS and pCS can be
compared directly; this compresses HA in panel B, whose accumulation is an order of
magnitude smaller ($\approx +1$\,\%).
}
\label{fig:dosing_comparison_bars}
\end{figure}
